%%%%%%%%%%%%%%%%%%%%%%%%%%%%%%%%%%%%%%%%%%%%%%%%%%%%%%%%%%%%%%%%%%%%%%%%%%%%%%%%

\pagestyle{empty}
\begin{abstract}
  Este Trabajo Fin de Grado aborda el diseño, desarrollo, despliegue y operación de una plataforma cloud para la gestión de transportes sobre AWS. El proyecto se plantea como un backend de alcance funcional acotado, pero con una carga técnica significativa en arquitectura, automatización, infraestructura como código, seguridad básica y validación operativa.

  La solución implementa una API REST en ASP.NET Core sobre PostgreSQL para gestionar transportes, vehículos, conductores, usuarios, asignaciones, transiciones de estado y eventos logísticos. La aplicación se empaqueta mediante Docker y se despliega en AWS usando Terraform, ECS Fargate, RDS PostgreSQL, ECR, Application Load Balancer, Route 53, ACM, Secrets Manager, SSM Parameter Store y CloudWatch. También se ha definido un flujo de despliegue manual desde GitHub Actions con autenticación OIDC, evitando claves AWS persistentes.

  Como resultado, se ha obtenido un entorno \texttt{dev} real recreable y validado en AWS, con servicio ECS estable, migraciones EF Core ejecutadas como tarea puntual de ECS, verificación de salud, bootstrap de administrador, login con JWT, logs JSON correlables en CloudWatch, dashboard operativo, alarmas Terraform para ALB/ECS/RDS/API y configuración DNS/HTTPS mediante Route 53 y ACM. Tras capturar la evidencia operativa, el entorno \texttt{dev} se destruyó con Terraform para evitar costes residuales, conservando únicamente recursos base de bajo coste como el dominio y la zona DNS.

  \vspace*{25pt}
  \begin{segundoresumo}
    This Bachelor's Thesis addresses the design, development, deployment, and operation of a cloud-based transport management platform on AWS. The project is implemented as a deliberately scoped backend, with a strong technical focus on architecture, automation, infrastructure as code, basic security, and operational validation.

    The solution provides a REST API built with ASP.NET Core and PostgreSQL to manage transports, vehicles, drivers, users, assignments, state transitions, and logistics events. The application is packaged with Docker and deployed to AWS using Terraform, ECS Fargate, RDS PostgreSQL, ECR, Application Load Balancer, Route 53, ACM, Secrets Manager, SSM Parameter Store, and CloudWatch. A manual GitHub Actions deployment flow with OIDC authentication has also been defined, avoiding persistent AWS access keys.

    The resulting \texttt{dev} environment is reproducible and has been validated on AWS with a stable ECS service, one-off EF Core migrations, health checks, first-admin bootstrap, JWT login, JSON structured logs with request correlation, a CloudWatch dashboard, Terraform-managed alarms for ALB/ECS/RDS/API, and DNS/HTTPS configuration through Route 53 and ACM. After collecting operational evidence, the \texttt{dev} stack was destroyed with Terraform to avoid residual cloud costs, leaving only low-cost foundational resources such as the domain and DNS zone.
  \end{segundoresumo}
\vspace*{25pt}
\begin{multicols}{2}
\begin{description}
\item [\palabraschaveprincipal:] \mbox{} \\[-20pt]
  \begin{itemize}
    \item Computación en la nube
    \item DevOps
    \item Infraestructura como código
    \item AWS
    \item Terraform
    \item ASP.NET Core
    \item Gestión de transportes
  \end{itemize}
\end{description}
\begin{description}
\item [\palabraschavesecundaria:] \mbox{} \\[-20pt]
  \begin{itemize}
    \item Cloud computing
    \item DevOps
    \item Infrastructure as code
    \item AWS
    \item Terraform
    \item ASP.NET Core
    \item Transport management
  \end{itemize}
\end{description}
\end{multicols}

\end{abstract}
\pagestyle{fancy}

%%%%%%%%%%%%%%%%%%%%%%%%%%%%%%%%%%%%%%%%%%%%%%%%%%%%%%%%%%%%%%%%%%%%%%%%%%%%%%%%
